%!TEX root = hems_proj.tex
%%%%%%%%%%%%%%%%%%%%%%%%%%%%%%%%%%%%%%%%%%%%%%%%%%%%%%%%%%%%%%%%%%%%%%%
\chapter{Irodalmi Áttekintés és Fogalmi Alapok}\label{ch:ALAPOK}
%%%%%%%%%%%%%%%%%%%%%%%%%%%%%%%%%%%%%%%%%%%%%%%%%%%%%%%%%%%%%%%%%%%%%%%

\begin{osszefoglal}
Áttekintjük a HEMS terület fő irányait, valamint bevezetünk néhány alapfogalmat a modell értelmezéséhez.
\end{osszefoglal}

\section{Irodalmi háttér}

A HEMS rendszerek optimalizálása tipikusan korlátozott, nemlineáris feladat, amelyben a helyi optimumok és a bizonytalan bemenetek miatt a metaheurisztikus megoldások széles körben elterjedtek \cite{Ismail2014}. A sztochasztikus termelés és a tárolás kezelési nehézségei a rendszer rugalmasságát és korlátkezelést helyezik előtérbe \cite{Fadaee2012}. A különböző algoritmusok jellemzően trade-offokkal dolgoznak a gyorsaság és a pontosság, illetve az átlagos és a legjobb megoldás minősége között, ezért az összehasonlító vizsgálat módszertani szempontból különösen indokolt.

\section{Optimalizálási modell és döntési változók}

A jelenlegi implementációban a napi optimalizálási feladat 24 órás horizonton fut. A döntési vektor két fő részre bontható: az akkumulátor teljesítmény-vektorra -- $p_{batt}(t)$, ahol $t=0,\dots,23$ -- és az eltolható eszközök indítási idejére, ahol $s_j$ jelöli a $j$-edik eszköz indulási időpontját.

Az energiamérleg óránként a következőképpen írható fel:
\begin{equation}
P_{grid}(t) = P_{load}(t) + P_{shift}(t) - P_{pv}(t) + p_{batt}(t)
\end{equation}
ahol $P_{shift}(t)$ az időzített fogyasztók -- például mosógép, mosogatógép -- által az adott órában hozzáadott terhelés.

A napi célfüggvény a hálózati energia költségét és a korlátsértési büntetéseket egyesíti:
\begin{equation}
J = \sum_{t=0}^{23} C_t(P_{grid}(t),\pi(t)) + \lambda_{soc} \cdot \Phi_{soc} + \lambda_{end} \cdot \Phi_{end}
\end{equation}
ahol $\pi(t)$ az adott órás ár, $\Phi_{soc}$ a SOC-határok sértésének büntetése, $\Phi_{end}$ pedig a napvégi SOC-visszatérés büntetése.

\section{Korlátok és büntetéses kezelés}

A modellben explicit korlátok szerepelnek mind az akkumulátorra, mind az időzített eszközökre vonatkozóan. Az akkumulátornál az állapotváltozó (SOC) minimális és maximális értéke közé szorítás történik, amelyet kvadratikus büntetéssel egészítünk ki. Ez a megközelítés a gyakorlatban stabilabb keresést ad, mint a tisztán „hard" eldobásos korlátkezelés, mivel a büntetőfüggvény folytonosan vezeti az optimalizálót a megvalósítható tartomány felé.

Az időzített eszközöknél az indítási idő folytonos döntési változóként jelenik meg, de a kiértékelés során egész órára kerekítés történik. Ez a modell egyszerre tartja meg a metaheurisztikák folyamatos keresési előnyeit, miközben fizikailag értelmezhető ütemezést eredményez.

\section{Vétel és visszatáplálás aszimmetriája}

A jelenlegi költségmodellben a hálózati vételezés teljes órás áron történik, míg a visszatáplálás csökkentett szorzóval -- jellemzően az ár 30\%-ával -- kerül elszámolásra. Ez az aszimmetria valós üzemeltetési helyzeteket közelít: a túltermelés nem azonos értékű a csúcsidős fogyasztás kiváltásával. Emiatt az optimalizáló természetes módon olyan stratégiákat preferál, amelyek inkább a lokális önfogyasztást és a jól időzített akkumulátor-használatot valósítják meg, semmint a felesleges visszatáplálást.